\documentclass[12pt,a4paper]{article}
\usepackage[utf8]{inputenc}
\usepackage[T1]{fontenc}
\usepackage[spanish]{babel}
\usepackage{amsmath,amssymb}
\usepackage{geometry}
\usepackage{booktabs}
\usepackage{longtable}
\usepackage{xcolor}
\usepackage{graphicx}
\geometry{margin=2.5cm}

\title{Cuadro Comparativo de Distancias de Seguridad\\Líneas de Transmisión}
\author{Líneas de Transmisión}
\date{8 de junio de 2026}

\begin{document}
\maketitle

\section*{Introducción}

Este documento presenta un cuadro comparativo de las distancias de seguridad
aplicables a líneas eléctricas de alta tensión, basado en dos fuentes:

\begin{itemize}
    \item \textbf{223-2008}: Normativa española actualizada (ITC-LAT 07).
          Contiene ecuaciones y consideraciones actualizadas para el cálculo
          de distancias mínimas.
    \item \textbf{Distancias OHTL}: Documento original con representaciones
          gráficas y dibujos ilustrativos de las distintas configuraciones
          de distancia.
\end{itemize}

\newpage

\section{Cuadro comparativo}

\begin{longtable}{p{3.5cm}p{5cm}p{5cm}}
\toprule
\textbf{Tipo de distancia} & \textbf{223-2008 (normativa)} & \textbf{OHTL (dibujos)} \\
\midrule
\endhead

\bottomrule
\endfoot

\textbf{Distancia al suelo}
&
Altura mínima de conductores sobre el terreno:
\begin{itemize}
    \item $5$ m en terrenos sin circulación rodada
    \item $6$ m en terrenos con circulación rodada
    \item $6$ m mínimo general para conductores recubiertos
    \item $1$ m sobre altura máxima de maquinaria (mín. $6$ m)
\end{itemize}
(Art. 6.4)
&
En OHTL se muestran esquemas de la distancia vertical
desde el conductor (con flecha máxima) hasta el nivel del
terreno, caminos y superficies de agua, ilustrando las
cotas mínimas según el tipo de terreno.
\\

\addlinespace

\textbf{Distancia entre conductores}
&
Para conductores recubiertos, distancia mínima:
\[
D = \sqrt{F^2 + L^2} + D_{pp} + K
\]
con mínimo absoluto de $0.2$ m.
(Art. 6.3)
&
OHTL presenta diagramas de la separación horizontal
y vertical entre conductores en el mismo vano,
considerando oscilación por viento y flecha máxima.
\\

\addlinespace

\textbf{Cruce con líneas AT}
&
Cruce con conductores desnudos:
\begin{itemize}
    \item Distancia vertical mínima: $0.5$ m ($\le 30$ kV)
    \item Para $>30$ kV: aplicar distancias ITC-LAT 07 \S 5.6.1
    \item Línea aislada siempre por debajo de la desnuda
\end{itemize}
(Art. 6.5.1)
&
OHTL ilustra la configuración de cruce entre líneas,
mostrando la disposición de apoyos y la distancia
vertical de separación entre conductores de ambas líneas.
\\

\addlinespace

\textbf{Paralelismo con líneas AT}
&
Separación mínima entre trazas: $1.5 \times$ altura del apoyo
más alto. Mínimo $0.5$ m para líneas $\le 30$ kV.
(Art. 6.5.2)
&
OHTL muestra en planta y perfil la disposición de
líneas paralelas, con las distancias horizontales entre
ejes y la zona de servidumbre.
\\

\addlinespace

\textbf{Cruce con BT o}
\textbf{telecomunicación}
&
Distancia vertical mínima:
\begin{itemize}
    \item Con BT: $0.5$ m (cables aislados), $1$ m (conductores recubiertos)
    \item Con telecomunicación: $1$ m (aislados), $1.5$ m (recubiertos)
\end{itemize}
(Art. 6.6)
&
OHTL detalla con dibujos la disposición de los cruces,
indicando la separación vertical y la ubicación relativa
de apoyos.
\\

\addlinespace

\textbf{Distancia a carreteras}
&
Cruce sobre carretera:
\begin{itemize}
    \item Distancia vertical mínima: $7$ m sobre la rasante
    \item Aplican prescripciones especiales del art. 6.2
\end{itemize}
(Art. 6.7)
&
OHTL representa gráficamente el cruce sobre calzada,
con la cota de $7$ m y la distancia a los bordes de la vía.
\\

\addlinespace

\textbf{Distancia a ferrocarriles}
&
\begin{itemize}
    \item Sin electrificar: $7$ m sobre cabezas de carriles
    \item Electrificados: $4$ m sobre conductor más alto
          (líneas de energía, telefónica y telegráfica)
\end{itemize}
(Arts. 6.8 y 6.9)
&
OHTL ilustra la configuración de cruce sobre vías
férreas, la altura sobre carril y las distancias laterales
a los postes de catenaria.
\\

\addlinespace

\textbf{Distancia a teleféricos}
&
Distancia vertical mínima: $5$ m entre conductor y
parte más elevada del teleférico. El teleférico debe
ponerse a tierra en dos puntos.
(Art. 6.10)
&
OHTL muestra esquemas del cruce con teleféricos,
con las distancias verticales y la ubicación de las
tomas de tierra.
\\

\addlinespace

\textbf{Distancia a ríos y canales}
&
Se aplica lo indicado en ITC-LAT 07 \S 5.11,
con prescripciones especiales del art. 6.2.
(Art. 6.11)
&
OHTL ilustra el cruce sobre cursos de agua navegables,
con la altura mínima sobre el nivel máximo de agua
y las distancias a los márgenes.
\\

\addlinespace

\textbf{Distancia a antenas y}
\textbf{pararrayos}
&
Distancia mínima:
\begin{itemize}
    \item $1$ m para cables aislados reunidos en haz
    \item $1.5$ m para conductores recubiertos
\end{itemize}
Prohibido fijar antenas en apoyos de líneas AT.
(Art. 6.12)
&
OHTL representa la separación mínima entre
conductores y antenas/pararrayos cercanos.
\\

\addlinespace

\textbf{Paso por bosques y}
\textbf{arbolado}
&
Distancia mínima desde conductor a masa de arbolado:
\begin{itemize}
    \item $2$ m para líneas de $30$ kV
    \item $1.5$ m para líneas $\le 20$ kV
\end{itemize}
(Art. 6.13.1)
&
OHTL detalla la franja de corte de arbolado y la
distancia de seguridad entre conductores y copas
de los árboles.
\\

\addlinespace

\textbf{Distancia a edificaciones}
&
Se aplica lo especificado en ITC-LAT 07 \S 5.12.2
para conductores recubiertos.
(Art. 6.13.2)
&
OHTL muestra las distancias horizontales y verticales
desde los conductores a fachadas, techos y otras partes
de las edificaciones.
\\

\end{longtable}

\newpage

\section*{Resumen de distancias mínimas}

\begin{center}
\begin{tabular}{lll}
\toprule
\textbf{Concepto} & \textbf{Cables aislados} & \textbf{Conductores recubiertos} \\
\midrule
Suelo (sin circulación)    & $5$ m                & $6$ m \\
Suelo (con circulación)    & $6$ m                & $6$ m \\
Carreteras                 & $7$ m                & $7$ m \\
Ferrocarril sin elec.      & $7$ m                & $7$ m \\
Ferrocarril electrificado  & $4$ m                & $4$ m \\
Cruce con BT               & $0.5$ m              & $1$ m \\
Cruce con telecomunicación & $1$ m                & $1.5$ m \\
Antenas/pararrayos         & $1$ m                & $1.5$ m \\
Teleféricos                & $5$ m                & $5$ m \\
Arbolado ($30$ kV)         & ---                  & $2$ m \\
Arbolado ($\le 20$ kV)     & ---                  & $1.5$ m \\
\midrule
\textbf{Fuente} & \multicolumn{2}{l}{223-2008 (ITC-LAT 07, cap. 6)} \\
\bottomrule
\end{tabular}
\end{center}

\vspace{1cm}

\noindent\textbf{Nota:} Para las representaciones gráficas detalladas de cada
configuración, consultar el documento \textit{Distancias OHTL}, el cual
contiene los dibujos ilustrativos de cada tipo de distancia de seguridad
aquí descrita.

\end{document}
